\section{Discusión}
\label{sec:discusión}

El ajuste conjunto reprodujo la biomasa de ambas dietas (RMSE de 24.5 y 12.1~mg frente a masas máximas de 140 a 220~mg; día 14 dentro del 9~\%) y el orden y la forma temporal de las tasas de gas, pero dejó al descubierto un desacuerdo estructural entre los dos bloques de observación: las cotas de los cierres saturados quedaron sin alcanzar en el caso base y en los nueve escenarios (sec.~\ref{sec:sensibilidad}) y el gas de D4 mostró un desfase temporal que ninguna recalibración corrigió ($R^2$ siempre por debajo de $-0.15$). El modelo no sostiene a la vez el pico respiratorio temprano que exigen las tasas medidas y la trayectoria de crecimiento observada: esa tensión, y no un punto de calibración particular, es la principal limitación del núcleo actual. Aun con ella, el trabajo pone a prueba el modelo DEB de \textit{H. illucens} ---validado originalmente solo contra crecimiento y asimilación \parencite{eriksenDynamicModellingFeed2022}--- con un tipo de datos nuevo, gas separado por fuente y crecimiento simultáneo, y la concordancia encontrada da soporte independiente, aunque parcial, a su ecuación de producción de CO$_2$ \parencite{eriksenDynamicModellingFeed2022,eriksenMetabolicPerformanceFeed2024}.

El cierre emergente del ciclo por instares resultó esencial y robusto: sin el decrecimiento de $B_{\max}(t)$ el modelo no reproduce el colapso de las emisiones, y el instante del cese ($t_{\mathrm{stop}} = 11.4$ y 12.5~d en D1 y D4) cayó dentro de las ventanas observadas y apenas se movió en los nueve escenarios (11.1 a 11.9~d y 11.5 a 12.8~d), incluso perturbando $t_{pm}$ en $\pm 20$~\%. Los datos determinan el cruce $B(t) = B_{\max}(t)$ y no a $t_{pm}$ y $\rho$ por separado (sec.~\ref{sec:sensibilidad}), y el bootstrap apunta en la misma dirección: $f$ y el cierre emergente son las cantidades mejor identificadas. El cese es así una propiedad emergente del crecimiento, no un interruptor impuesto, y el colapso se atribuye a la prepupa y no al agotamiento del sustrato: el lecho conservó cerca del 87~\% de su materia seca al final y el crecimiento cesó a pesar de la segunda carga del día~5. El suelo de mantenimiento en prepupa, anclado a la medición de larvas en ayuno (14.3~ppm/min), es coherente con un estadio no alimentario sostenido por reservas \parencite{arreseInsectFatBody2010,llandresDynamicEnergyBudget2015}. Operativamente, el colapso de la tasa de CO$_2$ marca el fin del ciclo productivo: una señal no invasiva para decidir la cosecha, cuya anticipación es el resultado más robusto del modelo.

El factor de calidad calibrado ($f = 0.802$ para D4, con D1 como referencia) indica que los residuos de frutas fermentados se asimilan al 80~\% de la tasa del alimento estándar, en línea con el desempeño dependiente de la dieta documentado para la especie \parencite{laganaroGrowthMetabolicPerformance2021a,eriksenMetabolicPerformanceFeed2024}: una reducción del 20~\% en la asimilación es un costo moderado frente al beneficio de valorizar residuos. La lectura de $f$ como calidad nutricional, y no como un simple parámetro de ajuste, quedó respaldada por la validación externa con \textcite{schonUsingGasExchange2026}: los factores recuperados decrecieron de forma monótona con el contenido de proteína de las cuatro dietas (0.907 a 0.733), sin reajustar ningún otro parámetro. La descomposición de fuentes, combinada con curvas de crecimiento, se muestra así como una vía para caracterizar la calidad de sustratos nuevos en modelos DEB.

El módulo microbiano calibró valores similares entre dietas ($k_{\mathrm{ref}} = 0.441$ y $0.386$~d$^{-1}$, $Y_{CO_2} = 1.47$ y $1.19$~g~g$^{-1}$ en D1 y D4), lo que sugiere una actividad microbiana parecida en ambos sustratos. La contribución del lecho fue menor frente al pico larvario (13 a 34~\% del total antes del colapso), pero domina al inicio y al final del ciclo (más del 80~\% al final): los estudios que reportan emisiones totales \parencite{ermolaevGreenhouseGasEmissions2019,parodiBioconversionEfficienciesGreenhouse2020a,boakye-yiadomGreenhouseGasEmissions2022,lindbergProcessEfficiencyGreenhouse2022a} mezclan ambas fuentes y la fracción microbiana es fuertemente dependiente del tiempo, lo que condiciona las comparaciones entre sustratos y días. La biomasa microbiana creciente reproduce la tendencia de los controles ($R^2 = 0.967$ y $0.983$; RMSE de 1.0 y 1.5~ppm/min) y respalda ese crecimiento como la explicación del aumento. El $Y_{CO_2}$ de D1 coincide con el máximo estequiométrico de la oxidación completa de carbohidratos ($\approx 1.47$~g~g$^{-1}$), consistente con la fracción proteica del alimento estándar (21.4~\% medida, Cuadro~\ref{tab:dietas}); el valor menor de D4 es consistente con un sustrato ya parcialmente fermentado, que perdió parte de su carbono antes del ensayo.

En unidades de ingeniería, las emisiones acumuladas del ciclo completo (0 a 15~d) ascienden a $72$ y $65$~g de CO$_2$ por kg de alimento fresco en D1 y D4 ($83$ y $80$~g~kg$^{-1}$ en base seca), con fracción microbiana acumulada del 26 y el 23~\%. Estos totales caen dentro del rango reportado para la bioconversión con \textit{H. illucens} ---96~g~kg$^{-1}$ de residuo alimentario en base húmeda \parencite{ermolaevGreenhouseGasEmissions2019} y 47 a 147~g~kg$^{-1}$ según el residuo y su pretratamiento \parencite{lindbergProcessEfficiencyGreenhouse2022a}--- y la fracción microbiana es comparable al 34~\% estimado por \textcite{parodiBioconversionEfficienciesGreenhouse2020a} a partir de su balance de masa. Dos salvedades: los totales son conservadores, porque el modelo subestima las emisiones tempranas, y al ser carbono biogénico este CO$_2$ no cuenta como emisión neta en los balances de calentamiento global \parencite{mertenatBlackSoldierFly2019}; su valor aquí es como trazador del destino del carbono y como base de la separación de fuentes.

Entre las limitaciones del experimento, los cierres hasta saturación restringieron el rango medible (dos puntos de D1 quedaron como cotas inferiores); cierres más cortos o un sensor de mayor rango mejorarían la estimación. Las curvas de crecimiento provienen de muestras de 10 larvas por recipiente: la masa media a granel al final del ensayo ($\approx 257$~mg por larva) supera con creces el último punto de las curvas muestreadas ($\approx 173$~mg en D1), lo que sugiere un sesgo de tamaño (las larvas grandes se capturan más fácil); como la calibración usó las curvas muestreadas, las masas absolutas del modelo heredan esa incertidumbre, aunque la trayectoria temporal y $t_{\mathrm{stop}}$ no deberían verse afectadas. Las métricas de bioconversión en base húmeda quedaron confundidas por la evaporación ---que el modelo no representa--- y no se usaron como validación cuantitativa.

La limitación conceptual principal es estructural: el modelo subestima sistemáticamente la fase explosiva del crecimiento (día 4, alrededor de un 32~\% en ambas dietas), porque el freno de Verhulst del último instar entra desde $B \approx B_{\max,0}/\delta_I \approx 10$~mg, y el dato pide más asimilación temprana de la que el rango de $a_{\max}$ permite, cuya cota superior se alcanzó en el caso base y en cuatro de los ocho escenarios perturbados. Estos rasgos apuntan a mecanismos ausentes del núcleo: un compartimento de reservas de movilización rápida que desacople la respiración del crecimiento estructural, una dependencia respiratoria de la ingesta reciente, o reservorios de sustrato de accesibilidad distinta para larvas y microbios que un solo compartimento $D_M$ bien mezclado no puede representar.

Finalmente, $k_{\mathrm{ref}}$ e $Y_{CO_2}$ entran al modelo como producto y están parcialmente correlacionados, y el número de larvas (una postura aproximada de 700) propaga su incertidumbre linealmente a las tasas absolutas.

La validación externa con \textcite{ermolaevGreenhouseGasEmissions2019} (sec.~\ref{sec:resultados_ermolaev}) respalda la transferibilidad del modelo: con la parametrización calibrada fija y reajustando solo los cuatro parámetros de dieta y cierre del ciclo ($f$, $t_{pm}$, $\rho$ y $B_{\max,0}$), se reprodujo un ensayo con otra dieta (residuos de alimentos), otro esquema de alimentación (pulsos discretos) y otra temperatura (28~$^{\circ}$C). El valor $f = 0.883$ queda entre los de este trabajo (1.0 y 0.802), coherente con un residuo de calidad intermedia, y los valores recuperados de $t_{pm}$ ($\approx 13$~d) y $B_{\max,0}$ ($\approx 95$~mg) son compatibles con la duración del ciclo y la cepa de ese estudio. Que el desacuerdo restante ---arranque subestimado y pico adelantado alrededor de un día y medio--- sea el mismo rasgo de fase temprana de los datos propios refuerza que la limitación es estructural y no de calibración; el cierre emergente, además, reproduce el final del ciclo sin sobreestimar la biomasa, sin requerir mecanismos adicionales.

La discrepancia de fase del gas de D4 admite además una lectura ecológica. La descomposición de fuentes asume que la actividad microbiana del tratamiento es la misma que la del control de alimento solo, pero las larvas interactúan con el lecho: lo remueven y airean, excretan nitrógeno lábil y pastan sobre su microbiota \parencite{liuBlackSoldierFly2008,zhangBlackSoldierFly2021}. En un sustrato ya fermentado como D4, con una comunidad microbiana establecida, ese estímulo \parencite{qinBlackSoldierFly2026} (y su consumo) puede generar un pulso microbiano temprano que la resta tratamiento menos control asigna íntegramente a la larva, adelantando el pico aparente de emisión sin alterar su magnitud ---exactamente el patrón observado. Se suman dos candidatos complementarios: un costo respiratorio de actividad o detoxificación asociado a los productos de la fermentación, que el núcleo DEB no representa, y la amortiguación del CO$_2$ disuelto en el lecho más húmedo. La hipótesis es consistente con que el modelo reproduzca D1, de comunidad microbiana inicial pobre e interacción débil. Discriminar entre estos mecanismos requiere datos que este experimento no registró ---biomasa microbiana del lecho con y sin larvas, o respirometría con lecho estéril frente a inoculado---: su cuantificación permitiría dotar al modelo de un término explícito de interacción larva--microbio, convirtiendo la limitación aquí declarada en la hipótesis central del siguiente estudio.

El protocolo de validación por descomposición de fuentes (tratamiento, control de alimento solo y ayuno), combinado con curvas de crecimiento no destructivas, es transferible a otros insectos de bioconversión \parencite{nielsenMetabolicPerformanceMealworms2025}, y la estructura del modelo admite extensiones futuras: la inclusión de O$_2$ y N$_2$O, la alimentación continua, la heterogeneidad espacial del lecho y el análisis de identificabilidad estructural de los parámetros del cierre del ciclo.
